\documentclass[12pt]{article}
\usepackage{geometry}
\geometry{margin=1in}

% ----------------------------------------------------------
%  Linux Penguin Theme — Based on Recommended Colors
% ----------------------------------------------------------

% --- COLOR SETUP ---
\usepackage{xcolor}

\definecolor{cream}{RGB}{246,242,219}
\definecolor{teamred}{RGB}{209,57,64}
\definecolor{navy}{RGB}{40,46,87}
\definecolor{softwhite}{RGB}{250,250,230}
\definecolor{accentyellow}{RGB}{240,200,60}

% Extra utility colors
\definecolor{muted}{RGB}{120,125,140}
\definecolor{darkcream}{RGB}{230,225,205}

% ----------------------------------------------------------
%  FONT SETUP (Optional but Recommended)
%  Works best with XeLaTeX or LuaLaTeX
% ----------------------------------------------------------
\usepackage{fontspec}

% Main document font (Google fonts — upload to Overleaf)
% If you don't upload fonts, comment these two lines.
\setmainfont{SourceSans3}[
  Path=./,
  ItalicFont  = *-Italic,
  BoldFont    = *-Bold
]

\newfontfamily\headingfont{ScienceGothic}[
  Path=./,
]

% Font Awesome for icons
\usepackage{fontawesome5}

% ----------------------------------------------------------
%  SECTION STYLE
% ----------------------------------------------------------
\usepackage{titlesec}

% Section
\titleformat{\section}
  {\Large\headingfont\bfseries\color{navy}}
  {\thesection}
  {1em}
  {\rule{\linewidth}{1pt}\vspace{0.5em}\\}

% Subsection
\titleformat{\subsection}
  {\large\headingfont\bfseries\color{teamred}}
  {\thesubsection}
  {0.8em}
  {}

% Subsubsection
\titleformat{\subsubsection}
  {\bfseries\color{navy}}
  {\thesubsubsection}
  {0.6em}
  {}

% ----------------------------------------------------------
%  PAGE STYLE
% ----------------------------------------------------------
\usepackage{fancyhdr}
\pagestyle{fancy}

\fancyhf{} % clear all

\lhead{\textcolor{navy}{\headingfont \courseTitle}}
\rhead{\textcolor{teamred}{\thepage}}
\renewcommand{\headrule}{\color{teamred}\hrule height 1pt}

% ----------------------------------------------------------
%  HYPERLINK STYLE
% ----------------------------------------------------------
\usepackage{hyperref}

\hypersetup{
  colorlinks = true,
  linkcolor  = teamred,
  urlcolor   = navy,
  citecolor  = accentyellow
}

% ----------------------------------------------------------
%  CUSTOM BOX ENVIRONMENTS
% ----------------------------------------------------------
\usepackage{tcolorbox}
\tcbuselibrary{skins,breakable}

% Info box
\newtcolorbox{infobox}{
  enhanced,
  breakable,
  colback   = softwhite,
  colframe  = navy,
  coltext   = navy,
  boxrule   = 1pt,
  arc       = 4pt,
  left=6pt, right=6pt, top=6pt, bottom=6pt
}

% Warning box
\newtcolorbox{warningbox}{
  enhanced,
  breakable,
  colback   = accentyellow!30,
  colframe  = accentyellow!80!navy,
  coltext   = navy,
  boxrule   = 1pt,
  arc       = 4pt,
  left=6pt, right=6pt, top=6pt, bottom=6pt
}

% Highlight box
\newtcolorbox{highlightbox}{
  enhanced,
  breakable,
  colback   = teamred!10,
  colframe  = teamred,
  coltext   = navy,
  boxrule   = 1pt,
  arc       = 4pt,
  left=6pt, right=6pt, top=6pt, bottom=6pt
}

% ----------------------------------------------------------
%  ITEMIZE / ENUMERATE STYLE
% ----------------------------------------------------------
\usepackage{enumitem}

\setlist[itemize]{
  label=\textcolor{teamred}{\large\textbullet},
  leftmargin=1.2em,
  itemsep=4pt
}

\setlist[enumerate]{
  label=\textcolor{navy}{\arabic*.},
  leftmargin=1.4em,
  itemsep=4pt
}

% ----------------------------------------------------------
%  TITLE STYLE
% ----------------------------------------------------------
\usepackage{titling}

\pretitle{
  \begin{center}
  \headingfont\bfseries\LARGE\color{navy}
}
\posttitle{
  \end{center}
}

\preauthor{
  \begin{center}
  \color{teamred}
}
\postauthor{
  \end{center}
}

% ----------------------------------------------------------
%  OPTIONAL — SOFT PAGE BACKGROUND
%  Comment out if you want white pages
% ----------------------------------------------------------
\usepackage{pagecolor}
\pagecolor{cream}
\color{navy}

% ----------------------------------------------------------
% CUSTOM COMMANDS FOR TEMPLATE
% ----------------------------------------------------------
\newcommand{\courseTitle}{Software Development on Linux}
\newcommand{\courseDate}{26-12-2025}
\newcommand{\authorName}{Khaled Mahfouz}
\newcommand{\contactInfo}{\texttt{@khaledmhfz}}

% Icon commands using Font Awesome
\newcommand{\devicon}{\faUserTie}
\newcommand{\mindseticon}{\faBrain}
\newcommand{\updateicon}{\faSync}
\newcommand{\terminalicon}{\faTerminal}
\newcommand{\toolicon}{\faTools}
\newcommand{\compileicon}{\faCogs}
\newcommand{\servericon}{\faServer}
\newcommand{\phpicon}{\faPhp}
\newcommand{\nodeicon}{\faNodeJs}
\newcommand{\reacticon}{\faReact}
\newcommand{\javaicon}{\faJava}
\newcommand{\browsericon}{\faFirefox}
\newcommand{\codeicon}{\faCode}
\newcommand{\arduinoicon}{\faMicrochip}
\newcommand{\debugicon}{\faBug}
\newcommand{\makeicon}{\faFileCode}
\newcommand{\pythonicon}{\faPython}
\newcommand{\appicon}{\faDesktop}
\newcommand{\repoicon}{\faBox}
\newcommand{\warningicon}{\faExclamationTriangle}
\newcommand{\dockericon}{\faDocker}
\newcommand{\rocketicon}{\faRocket}
\newcommand{\bookicon}{\faBook}
\newcommand{\fireicon}{\faFire}
\newcommand{\computericon}{\faDesktop}

% ----------------------------------------------------------
% END OF THEME
% ----------------------------------------------------------

\begin{document}

% Title with minimal space
\begin{center}
\headingfont\Huge\bfseries\color{navy}
\codeicon\ \courseTitle\\[0.5em]
\Large\color{teamred}
\courseDate
\end{center}

\vspace{1em}
\begin{center}
\color{teamred}\rule{0.8\textwidth}{1pt}
\end{center}

\section*{0. Before We Start (\mindseticon\ Mindset \& Setup)}

\begin{itemize}
    \item Everything is a tool
    \item The terminal is your superpower
    \item You \textit{own} your system
\end{itemize}

Update first. Always.

\begin{highlightbox}
\texttt{\$ sudo apt update \&\& sudo apt upgrade -y}
\end{highlightbox}

Install the essentials:

\begin{highlightbox}
\texttt{\$ sudo apt install -y build-essential curl wget git software-properties-common}
\end{highlightbox}

\begin{infobox}
\bookicon\ \textbf{Docs:}
\begin{itemize}
    \item \href{https://wiki.debian.org/Apt}{\texttt{https://wiki.debian.org/Apt}}
    \item \href{https://help.ubuntu.com/community/AptGet/Howto}{\texttt{https://help.ubuntu.com/community/AptGet/Howto}}
\end{itemize}
\end{infobox}

\vspace{1em}
\begin{center}
\color{teamred}\rule{0.8\textwidth}{1pt}
\end{center}

\section*{1. Compiling Suckless Tools (\compileicon\ st, dmenu)}

Suckless tools are \textbf{minimal, fast, and opinionated}. You compile them yourself. No training wheels.

\subsection*{Install dependencies}

\begin{highlightbox}
\texttt{\$ sudo apt install -y libx11-dev libxft-dev libxinerama-dev}
\end{highlightbox}

\subsection*{Clone and build \texttt{st} (simple terminal)}

\begin{highlightbox}
\texttt{\$ git clone https://git.suckless.org/st}\\
\texttt{\$ cd st}\\
\texttt{\$ sudo make install clean}
\end{highlightbox}

Same idea for \texttt{dmenu}:

\begin{highlightbox}
\texttt{\$ git clone https://git.suckless.org/dmenu}\\
\texttt{\$ cd dmenu}\\
\texttt{\$ sudo make install clean}
\end{highlightbox}

\begin{infobox}
\fireicon\ \textbf{Pro tip:} You \textit{edit config.h and recompile}. That's the suckless way.
\end{infobox}

\begin{infobox}
\bookicon\ \textbf{Official docs:}
\begin{itemize}
    \item \href{https://suckless.org/st/}{\texttt{https://suckless.org/st/}}
    \item \href{https://suckless.org/dmenu/}{\texttt{https://suckless.org/dmenu/}}
\end{itemize}
\end{infobox}

\vspace{1em}
\begin{center}
\color{teamred}\rule{0.8\textwidth}{1pt}
\end{center}

\section*{2. Run a Simple Server \& Share Files (\servericon)}

\subsection*{Python HTTP server (the classic)}

\begin{highlightbox}
\texttt{\$ python3 -m http.server 8000}
\end{highlightbox}

Now anyone on your network can access:

\begin{highlightbox}
\texttt{http://YOUR\_IP:8000}
\end{highlightbox}

\begin{warningbox}
\warningicon\ Warning: http.server is not recommended for production. It only implements basic security checks.
\end{warningbox}

\subsection*{Node.js quick server}

\begin{highlightbox}
\texttt{\$ npx serve .}
\end{highlightbox}

\begin{infobox}
\bookicon\ \textbf{Docs:}
\begin{itemize}
    \item \href{https://docs.python.org/3/library/http.server.html}{\texttt{https://docs.python.org/3/library/http.server.html}}
    \item \href{https://www.npmjs.com/package/serve}{\texttt{https://www.npmjs.com/package/serve}}
\end{itemize}
\end{infobox}

\vspace{1em}
\begin{center}
\color{teamred}\rule{0.8\textwidth}{1pt}
\end{center}

\section*{3. PHP \& XAMPP (\phpicon\ Yes, PHP Still Exists \faSmile)}

\begin{infobox}
\computericon\ Choose \textbf{One} of the following ways below
\end{infobox}

\subsection*{Install PHP}

\begin{highlightbox}
\texttt{\$ sudo apt install -y php php-cli php-mysql}
\end{highlightbox}

Run a PHP dev server:

\begin{highlightbox}
\texttt{\$ php -S localhost:8000}
\end{highlightbox}

\subsection*{XAMPP (All-in-One)}

Download from Apache Friends:

\begin{highlightbox}
\texttt{https://www.apachefriends.org/index.html}
\end{highlightbox}

Install:

\begin{highlightbox}
\texttt{\$ chmod +x xampp-linux-*-installer.run}\\
\texttt{\$ sudo ./xampp-linux-*-installer.run}
\end{highlightbox}

Start it:

\begin{highlightbox}
\texttt{\$ sudo /opt/lampp/lampp start}
\end{highlightbox}

\begin{infobox}
\bookicon\ \textbf{Docs:}
\begin{itemize}
    \item \href{https://www.php.net/manual/en/}{\texttt{https://www.php.net/manual/en/}}
    \item \href{https://www.apachefriends.org/faq_linux.html}{\texttt{https://www.apachefriends.org/faq\_linux.html}}
\end{itemize}
\end{infobox}

\vspace{1em}
\begin{center}
\color{teamred}\rule{0.8\textwidth}{1pt}
\end{center}

\section*{4. Node.js \& Your First React App (\reacticon)}

\subsection*{Install Node.js (LTS)}

\begin{highlightbox}
\texttt{\# Download and install nvm:}\\
\texttt{\$ curl -o- https://raw.githubusercontent.com/nvm-sh/nvm/v0.40.3/install.sh | bash}\\
\\
\texttt{\# in lieu of restarting the shell}\\
\texttt{\$ \textbackslash."\$HOME/.nvm/nvm.sh"}\\
\\
\texttt{\# Download and install Node.js:}\\
\texttt{\$ nvm install 24}
\end{highlightbox}

Verify:

\begin{highlightbox}
\texttt{\# Verify the Node.js version:}\\
\texttt{\$ node -v \# Should print "v24.12.0".}\\
\\
\texttt{\# Verify npm version:}\\
\texttt{\$ npm -v \# Should print "11.6.2".}
\end{highlightbox}

\subsection*{Create a React app}

\begin{highlightbox}
\texttt{\$ npx create-react-app my-app}\\
\texttt{\$ cd my-app}\\
\texttt{\$ npm start}
\end{highlightbox}

Boom. Browser opens. Magic.

\begin{infobox}
\bookicon\ \textbf{Docs:}
\begin{itemize}
    \item \href{https://nodejs.org/en/docs}{\texttt{https://nodejs.org/en/docs}}
    \item \href{https://react.dev/learn}{\texttt{https://react.dev/learn}}
\end{itemize}
\end{infobox}

\vspace{1em}
\begin{center}
\color{teamred}\rule{0.8\textwidth}{1pt}
\end{center}

\section*{5. Java Basics on Linux (\javaicon)}

\subsection*{Install OpenJDK}

\begin{highlightbox}
\texttt{\$ sudo apt install default-jdk}
\end{highlightbox}

You can also install specific versions of OpenJDK (e.g., Java 11 or Java 21) by searching and choosing the one you need:

\begin{highlightbox}
\texttt{\$ apt search openjdk}\\
\texttt{\$ sudo apt install openjdk-11-jdk}
\end{highlightbox}

Verify:

\begin{highlightbox}
\texttt{\$ java -version}\\
\texttt{\$ javac -version}
\end{highlightbox}

\subsection*{Hello World (because tradition)}

\begin{highlightbox}
\begin{verbatim}
public class Hello {
    public static void main(String[] args) {
        System.out.println("Hello Linux");
    }
}
\end{verbatim}
\end{highlightbox}

Compile \& run:

\begin{highlightbox}
\texttt{\$ javac Hello.java}\\
\texttt{\$ java Hello}
\end{highlightbox}

\begin{infobox}
\bookicon\ \textbf{Docs:}
\begin{itemize}
    \item \href{https://openjdk.org/}{\texttt{https://openjdk.org/}}
    \item \href{https://docs.oracle.com/javase/tutorial/}{\texttt{https://docs.oracle.com/javase/tutorial/}}
\end{itemize}
\end{infobox}

\vspace{1em}
\begin{center}
\color{teamred}\rule{0.8\textwidth}{1pt}
\end{center}

\section*{6. Browsers, VS Code \& Arduino IDE (\browsericon)}

\subsection*{Browsers}

\begin{highlightbox}
\texttt{\$ sudo apt install -y firefox}
\end{highlightbox}

Chrome:

\begin{highlightbox}
\texttt{https://www.google.com/chrome/}
\end{highlightbox}

\subsection*{VS Code}

Download from VS Code official website :

\begin{highlightbox}
\texttt{https://code.visualstudio.com/}
\end{highlightbox}

\subsection*{Arduino IDE}

\begin{highlightbox}
\texttt{\$ sudo apt install -y arduino}
\end{highlightbox}

\begin{infobox}
\bookicon\ \textbf{Docs:}
\begin{itemize}
    \item \href{https://code.visualstudio.com/docs/setup/linux}{\texttt{https://code.visualstudio.com/docs/setup/linux}}
    \item \href{https://docs.arduino.cc/}{\texttt{https://docs.arduino.cc/}}
\end{itemize}
\end{infobox}

\vspace{1em}
\begin{center}
\color{teamred}\rule{0.8\textwidth}{1pt}
\end{center}

\section*{7. Compilers \& Debugging Tools (\debugicon)}

\subsection*{C/C++}

\begin{highlightbox}
\texttt{\$ sudo apt install -y gcc g++ gdb}
\end{highlightbox}

Compile:

\begin{highlightbox}
\texttt{\$ gcc main.c -o main}
\end{highlightbox}

Debug:

\begin{highlightbox}
\texttt{\$ gdb ./main}
\end{highlightbox}

\begin{infobox}
\bookicon\ \textbf{Docs:}
\begin{itemize}
    \item \href{https://gcc.gnu.org/onlinedocs/}{\texttt{https://gcc.gnu.org/onlinedocs/}}
    \item \href{https://www.gnu.org/software/gdb/documentation/}{\texttt{https://www.gnu.org/software/gdb/documentation/}}
\end{itemize}
\end{infobox}

\vspace{1em}
\begin{center}
\color{teamred}\rule{0.8\textwidth}{1pt}
\end{center}

\section*{8. Makefiles \& AVR Toolchain (\makeicon)}

\subsection*{Makefile basics}

\begin{highlightbox}
\begin{verbatim}
all:
    gcc main.c -o main
\end{verbatim}
\end{highlightbox}

Run:

\begin{highlightbox}
\texttt{\$ make}
\end{highlightbox}

\subsection*{AVR toolchain}

\begin{highlightbox}
\texttt{\$ sudo apt install -y avr-libc avrdude gcc-avr}
\end{highlightbox}

\subsection*{Helper Makefile to compile AVR code}

\begin{highlightbox}
\begin{verbatim}
#MCU name
MCU = atmega16

# CPU frequency
F_CPU = 8000000UL

# Compiler and tools
CC = avr-gcc
OBJCOPY = avr-objcopy
SIZE = avr-size
DIR = build

# Source file
SRC = main.c

# Output files
TARGET = main
ELF = $(DIR)/$(TARGET).elf
HEX = $(DIR)/$(TARGET).hex
# Compiler flags
CFLAGS = -mmcu=$(MCU) -DF_CPU=$(F_CPU) -Os

hex: $(HEX)
elf: $(ELF)

$(ELF): $(SRC)
    mkdir -p $(DIR)
    $(CC) $(CFLAGS) -o $@ $^

$(HEX): $(ELF)
    $(OBJCOPY) -O ihex -R .eeprom $< $@
    $(SIZE) $<

clean:
    rm -rf build/

.PHONY: hex elf clean
\end{verbatim}
\end{highlightbox}

Used for Arduino \& AVR microcontrollers.

Run:

\begin{highlightbox}
\texttt{\$ make elf \# make elf binary}\\
\texttt{\$ make hex \# make hex binary}\\
\texttt{\$ make clean \# remove build/ directory}
\end{highlightbox}

\begin{infobox}
\bookicon\ \textbf{Docs:}
\begin{itemize}
    \item \href{https://www.gnu.org/software/make/manual/}{\texttt{https://www.gnu.org/software/make/manual/}}
    \item \href{https://www.nongnu.org/avr-libc/}{\texttt{https://www.nongnu.org/avr-libc/}}
\end{itemize}
\end{infobox}

\vspace{1em}
\begin{center}
\color{teamred}\rule{0.8\textwidth}{1pt}
\end{center}

\section*{9. Python, pip \& virtual environments (\pythonicon)}

\subsection*{Install Python}

\begin{highlightbox}
\texttt{\$ sudo apt install -y python3 python3-pip python3-venv}
\end{highlightbox}

\subsection*{Virtual environment (DO THIS)}

\begin{highlightbox}
\texttt{\$ python3 -m venv venv}\\
\texttt{\$ source venv/bin/activate}
\end{highlightbox}

Install packages:

\begin{highlightbox}
\texttt{\$ pip install requests bs4}
\end{highlightbox}

\begin{infobox}
\bookicon\ \textbf{Docs:}
\begin{itemize}
    \item \href{https://docs.python.org/3/tutorial/venv.html}{\texttt{https://docs.python.org/3/tutorial/venv.html}}
    \item \href{https://pip.pypa.io/en/stable/}{\texttt{https://pip.pypa.io/en/stable/}}
\end{itemize}
\end{infobox}

\vspace{1em}
\begin{center}
\color{teamred}\rule{0.8\textwidth}{1pt}
\end{center}

\section*{10. AppImage \& Desktop Shortcuts (\appicon)}

Run AppImage:

\begin{highlightbox}
\texttt{\$ chmod +x MyApp.AppImage}\\
\texttt{\$ ./MyApp.AppImage}
\end{highlightbox}

\subsection*{Desktop shortcut}

Create:

\begin{highlightbox}
\texttt{\$ \textasciitilde/.local/share/applications/myapp.desktop}
\end{highlightbox}

\begin{highlightbox}
\begin{verbatim}
[Desktop Entry]
Name=My App
Exec=/path/MyApp.AppImage
Type=Application
Icon=/path/icon.png
\end{verbatim}
\end{highlightbox}

\begin{infobox}
\bookicon\ \textbf{Docs:}
\begin{itemize}
    \item \href{https://docs.appimage.org/}{\texttt{https://docs.appimage.org/}}
\end{itemize}
\end{infobox}

\begin{infobox}
\computericon\ \textbf{Helping Tutorial:}
\begin{itemize}
    \item \href{https://mmbesar.github.io/Tutorials/appimages/}{\texttt{mmbesar}}
\end{itemize}
\end{infobox}

\vspace{1em}
\begin{center}
\color{teamred}\rule{0.8\textwidth}{1pt}
\end{center}

\section*{11. Update System \& Change Repositories (\repoicon)}

Edit sources:

\begin{highlightbox}
\texttt{\$ sudo nano /etc/apt/sources.list}
\end{highlightbox}

Update:

\begin{highlightbox}
\texttt{\$ sudo apt update \&\& sudo apt upgrade}
\end{highlightbox}

\begin{infobox}
\bookicon\ \textbf{Docs:}
\begin{itemize}
    \item \href{https://wiki.debian.org/SourcesList}{\texttt{https://wiki.debian.org/SourcesList}}
    \item \href{https://help.ubuntu.com/community/Repositories}{\texttt{https://help.ubuntu.com/community/Repositories}}
\end{itemize}
\end{infobox}

\vspace{1em}
\begin{center}
\color{teamred}\rule{0.8\textwidth}{1pt}
\end{center}

\section*{12. Kali Repositories (\warningicon\ Read This)}

\begin{warningbox}
\warningicon\ \textbf{Don't mix Kali repos with Ubuntu casually.}\\
You \textit{will} break things.
\end{warningbox}

Safe use case: \textbf{specific tools only}.

\begin{highlightbox}
\texttt{deb http://http.kali.org/kali kali-rolling main non-free contrib}
\end{highlightbox}

Use pinning if you \textit{must}.

\begin{infobox}
\bookicon\ \textbf{Docs:}
\begin{itemize}
    \item \href{https://www.kali.org/docs/general-use/kali-linux-sources-list-repositories/}{\texttt{kali-linux-sources-list-repositories [official]}}
\end{itemize}
\end{infobox}

\begin{infobox}
\computericon\ \textbf{Helping Tutorial:}
\begin{itemize}
    \item \href{https://youtu.be/YI1Q3R0TEYs?si=TO3KfOz4nq0sSIVw}{\texttt{Adding Kali linux repos to debian distros}}
\end{itemize}
\end{infobox}

\vspace{1em}
\begin{center}
\color{teamred}\rule{0.8\textwidth}{1pt}
\end{center}

\section*{13. Distrobox (Run Any Distro Safely)}

\begin{highlightbox}
\texttt{\$ sudo apt install -y podman}
\end{highlightbox}

Create container:

\begin{highlightbox}
\texttt{\$ distrobox create -n arch --image archlinux:latest}\\
\texttt{\$ distrobox enter arch}
\end{highlightbox}

Now you're in another distro. No VM. No pain.

\begin{infobox}
\bookicon\ \textbf{Docs:}
\begin{itemize}
    \item \href{https://distrobox.it/}{\texttt{https://distrobox.it/}}
\end{itemize}
\end{infobox}

\vspace{1em}
\begin{center}
\color{teamred}\rule{0.8\textwidth}{1pt}
\end{center}

\section*{14. Docker \& Containers \dockericon}

\subsection*{Install Docker}

\begin{highlightbox}
\texttt{\$ sudo apt install -y docker.io}\\
\texttt{\$ sudo usermod -aG docker \$USER}\\
\texttt{\$ newgrp docker}
\end{highlightbox}

Test:

\begin{highlightbox}
\texttt{\$ docker run hello-world}
\end{highlightbox}

Run a container:

\begin{highlightbox}
\texttt{\$ docker run -it ubuntu bash}
\end{highlightbox}

\begin{infobox}
\bookicon\ \textbf{Docs:}
\begin{itemize}
    \item \href{https://docs.docker.com/engine/install/ubuntu/}{\texttt{https://docs.docker.com/engine/install/ubuntu/}}
\end{itemize}
\end{infobox}

\begin{infobox}
\computericon\ \textbf{Helping Tutorial:}
\begin{itemize}
    \item \href{https://mmbesar.github.io/Tutorials/HS-Docker/}{\texttt{Home Server - Install and Use Docker}}
\end{itemize}
\end{infobox}

\vspace{1em}
\begin{center}
\color{teamred}\rule{0.8\textwidth}{1pt}
\end{center}

% Updated Footer
\vspace{\fill}
\begin{center}
\color{muted}
All rights reserved for \texttt{TogetherTeam} [NOT For Commercial Use] !!\\
contact the author via email \href{mailto:khaledmhfz2004@gmail.com}{\texttt{khaledmhfz2004@gmail.com}}
\end{center}

\end{document}